\section{讨论}
\label{sec:discussion}
\subsection{阴性结果说明了什么：正常传播与跨次记忆是不同能力}
M0 已经能够在 JO-CE 输入下产生正常回路响应，但在本轮两个间隔中没有逐次递减。这把“连线大致正确、能够传播信号”与“能在相应时间尺度保存刺激历史”区分开来。对初学者而言，关键不是模型有多少神经元，而是哪些状态能在两次刺激之间持续影响输出。

这一结果不能证明原模型在所有刺激、初态或参数下都没有慢动力学；也不能证明其他固定权重循环网络无法习惯化。我们没有扫描完整输入空间、系统搜索慢模态或改变原始参数。合理表述应始终附带条件：\emph{在所测试协议和参数范围内}，M0 不足以产生观察目标。

\subsection{充分性不是必要性，“最小扩展”不是唯一机制}
M1 表明，仅改变一组感觉输入的中央输出效能，就可以在固定图谱上得到所测的部分特征。这是一个模型内的\emph{充分性构造}：存在这样一组状态方程和参数，使现象出现。它不是\emph{必要性证明}，因为没有排除神经元适应、抑制增强、去抑制调节或其他慢状态机制。

$U=0$ 工程对照证明替换连接实现本身没有制造基线差异；它不等同于在真实动物中选择性阻断了某种分子机制，也不等于比较了所有候选记忆位置。M1 的作用点是 JO-CE 在脑内的兴奋性输出，而非经验证的习惯化特异性突触集合。只有进一步进行位置受约束的竞争模型与干预，才能回答机制在哪里、为什么在那里。

此外，M1 的强递减在一定程度上符合资源耗竭规则的直观预期。\textbf{“按会耗竭的规则建模，然后看到耗竭样下降”本身不足以构成新颖机制发现。}研究价值来自明确约束、阴性基线、实现对照，以及未来对未用于选择模型的协议和干预作出可被否定的预测。

\subsection{真实连接组的贡献尚未得到证明}
本轮保留了完整解剖结构，但没有比较裁剪子回路、随机图谱或低维模型。因此，不能把“在真实连接组上成功”改写成“只有真实连接组才能成功”。最小模型文献已经表明，小系统也可以展示下降与恢复\citep{smart2024}。

若要论证连接组有独特价值，应控制参数数量、拟合数据量和测试任务，并要求结构模型对留出刺激或细胞干预作出更好的预测。例如，两种模型都能拟合同一条 aBN1 曲线，但对抑制某个中间细胞的后果不同；若真实干预支持连接组模型，结构信息的贡献才得到更有力的证明。仅比较图的大小或视觉上更像大脑，不是有效对照。

\subsection{为什么恢复测试比单条下降曲线更有区分力}
短时耗竭、慢抑制和细胞适应都可能产生递减，但它们对休息长度、训练频率、追加刺激和旁路输入的预测可能不同。恢复测试因此有机会把“都能画出下降曲线”的模型分开。

本轮只有两个探测时间，且主读出是小整数计数，恢复曲线既稀疏又受量化影响。对于首次只产生少量脉冲的下游读出，少一个脉冲就可能导致很大的相对变化。这些因素使精确恢复速率估计不可靠。

还需避免把文献的数学结论直接移植到不同系统。Smart 等人的单状态最小模型具有明确的 H4(b) 限制\citep{smart2024}，但 M1 包含多个局部资源状态以及一个非线性全脑循环网络，并不是同一个数学对象。相同的局部恢复时间常数也不保证神经输出按同一个指数恢复。因此，本轮不把 H4(b) 自动判为通过或失败。

\subsection{正常功能损失与行为可解释性是实质性限制}
M1 的首次 aBN1 与下行响应降低，意味着“制造记忆”与“削弱整体传递”混在一起。一个好的生物候选模型不应仅在末期输出变小，还应在未训练时保持正常的刺激--输出关系，并在相关条件下保留应对新刺激的能力。

后续应预先设定正常功能门槛，例如要求多个下游读出在首次刺激时远离零计数地板；用直接激活下游或替代输入检验执行通路；同时观测输入细胞、中间层和输出层。若某个候选机制必须先把正常传递几乎关掉才能获得递减，就应把它视为重要失败模式，而不是只展示归一化后的漂亮曲线。

\subsection{本研究的主要局限}
\begin{enumerate}
\item \textbf{神经代理而非真实行为。}没有身体、运动反馈及脉冲到行为的校准。aBN1 计数不是梳理概率，aDN 发放也不等于动作执行。
\item \textbf{只测试一个输入强度和两种间隔。}没有覆盖弱/强刺激、更多时间尺度、其他模态及多个初态。
\item \textbf{同一图谱上的有限输入实现。}五个种子不能代表五个动物、基因型、个体连接差异或生物总体。
\item \textbf{冻结回放是一种控制而非生态真实性。}实际刺激往往逐次波动；本轮优先排除输入噪声混淆，尚未检验新的随机脉冲每次重抽时的稳健性。
\item \textbf{参数不是生理拟合结果。}$U$ 与恢复时间常数没有从 JO-CE 重复刺激数据估计，也未做敏感性或可辨识性分析。
\item \textbf{机制竞争与结构对照缺失。}没有证明 STD 优于抑制可塑性，也没有证明完整图谱优于小模型。
\item \textbf{关键习惯化特征缺失。}特异性、去习惯化、长期习惯化以及完整频率依赖恢复尚未确立。
\item \textbf{发表新颖性尚待审查。}本轮完成的是可重复计算先导研究，未进行穷尽性的新颖性检索，也没有足够证据作为完整生物机制论文直接投稿。
\end{enumerate}

\section{下一阶段：怎样把先导结果变成可发表的机制研究}
\label{sec:next}
以下为未来建议，不是本报告已经执行的实验。优先推进能区分假说的测试，而不是扩大参数搜索直到所有曲线都符合预期。

\subsection{竞争模型与判别实验}
建议至少比较三个候选：局部短时资源耗竭、中央抑制增强、抑制增强加解除抑制。各模型先通过相同的正常功能门槛，再用一部分协议约束少量共享参数；用从未参与拟合的协议评估预测。

\begin{table}[tbp]
\centering\small
\caption{建议的判别实验。预测须在新实验前由具体模型产生；下表不是已经测得的神经或行为结果。}
\label{tab:future}
\begin{tabularx}{\linewidth}{>{\raggedright\arraybackslash}p{3.2cm}YY}
\toprule
实验问题 & 设计要点 & 能排除或区分什么\\
\midrule
恢复有几种有效尺度？ & 从同一末态独立探测 0.2、0.5、1、2、5、10 s；检查平台与末态深度。 & 一种资源恢复尺度与多状态/多尺度模型的留出预测。\\
只是输出变弱了吗？ & 训练后直接驱动下游；检查未训练时的输出动态范围。 & 一般失能、低计数地板与保留潜在响应能力。\\
是否对刺激有选择？ & 选择能驱动可比较读出的 B，测试训练 A 对 A/B 的不同影响。 & 广泛疲劳与通路相关状态；不能使用无有效基线的 B。\\
B 是否解除 A 的递减？ & A--B--A、A--等时休息--A、未训练--B--A。 & 去习惯化、自然恢复与一般敏化。\\
抑制增强是否必要？ & 在明确标注的候选抑制细胞上设计选择性干预。 & 不同机制对旁路、沉默或增益改变的分歧。\\
真实结构是否有用？ & 同等数据预算比较最小模型、局部图、匹配随机图与全图。 & 结构带来的干预预测增益，而非仅有更高拟合自由度。\\
\bottomrule
\end{tabularx}
\end{table}

\subsection{什么时候应该转向嗅觉或视觉}
如果主问题转向刺激特异性和抑制性可塑性，嗅觉具有更直接的文献机制依据，但需补齐细胞标注、输入编码及分钟级时间尺度，不能把 JO-CE 参数直接照搬\citep{das2011}。如果主要目标是经典行为 hallmarks 与个体差异，light-off 跳跃更有成熟行为基础，但必须匹配视觉编码、相关输出及腹神经索/身体范围\citep{engel2009,boon2026}。

已有糖刺激基线可以成为跨模态工程对照，但唇瓣糖输入不能直接等同于足部糖的 PER 习惯化协议\citep{trisal2022}。选择模态应由要识别的机制和可获得的实验证据决定，而不是由最先跑通的演示锁定。

\section{结论}
在原始图谱、静态参数与本轮 JO-CE 协议下，固定连接全脑 LIF 没有出现跨次响应递减。仅对一组兴奋性输入输出边加入可恢复资源状态，即可在 aBN1 生成强递减、休息恢复和间隔依赖：短间隔的平均早晚期递减为 90.22\%，长间隔为 9.88\%。这些结果经过完整事件回读和输入配对校验。

最重要的结论不是“已经重建了会学习的果蝇”，而是：\textbf{正常感觉传递与跨次慢记忆必须分开检验；产生部分习惯化样现象的条件性充分性，并不等于真实机制识别或连接组必要性。}下一步应聚焦竞争机制、正常功能保留和留出干预预测。

\section*{数据、代码与研究透明性声明}
\addcontentsline{toc}{section}{数据、代码与研究透明性声明}
\textbf{代码与数据可得性。}本地项目包含固定协议、连续状态入口、独立分析器、逐种子结果、原始事件、输入实现及源码快照。已有公开仓库地址为：
\begin{center}\small\url{https://github.com/chenmzh/digital-fly-brain}\end{center}
本轮习惯化目录和本报告尚未推送；不能据该链接声称当前全部材料已公开。原始事件目前仅保留在本机，公开归档与长期数据托管仍待完成。

\textbf{伦理范围。}本轮只分析公开来源的处理后连接数据并进行计算仿真，没有新增活体动物或人体实验。原始数据的来源、使用条件与论文归属仍应遵守上游约定；这里不虚构伦理审批编号。

\textbf{署名、资助与利益冲突。}责任作者、单位、贡献分工、资助及利益冲突声明尚待人工确认，本文不据此虚构“无利益冲突”或任何资助信息。

\textbf{AI 辅助。}本项目的代码整理、文献归纳、分析辅助、制图和报告写作使用了 AI 编程助手。报告引用的数值来自实际运行和文件校验，但是否适合投稿、作者责任以及科学解释仍需人工审阅。
